\section{Restitución radial orientada y transporte simpléctico}
\subsection{Recuperación de Planck con orientación conservada}

Con los semiejes etiquetados de la elipse, definimos
\(R_\star=\sqrt{b_S/a_S}\). Entonces
\(q=R_\star^4=e^{-2\eta_{\rm el}}=(1-\xi)/(1+\xi)\),
donde \(\xi=C^*/A\) y \(\eta_{\rm el}=\operatorname{artanh}\xi\).
Esta razón radial conserva las marcas de ambos ejes.

\begin{theorem}[Recuperador radial de la sección de retorno]
\label{rec:thm:accion}
En la coordenatización positiva, si \(q=R_\star^4\), entonces
\begin{align}
 \eta_{\rm ret}&=\alpha\frac{499-501q}{1+q},\label{rec:eq:eta-radio}\\
 \hbar_{\rm ret}&=\alpha\frac{499-501q}{1+q}\,10^{-34}\mathcal U_S,
 \qquad\hbar_{\rm pre}=H_5\,10^{-34}\mathcal U_S,\label{rec:eq:planck-radio}\\
 R_{\rm act}&=\frac{\hbar_{\rm pre}}{\hbar_{\rm ret}}
 =\frac{H_5(1+q)}{\alpha(499-501q)}.\label{rec:eq:ratio-radio}
\end{align}
La rama \(\eta_{\rm ret}>0\) equivale a \(q<499/501\), dentro de \(q>0\).
\end{theorem}
\begin{proof}
Las identidades angulares dan \(\eta_{\rm ret}=\alpha(500\xi-1)\). Se sustituye \(\xi=(1-q)/(1+q)\); el numerador es \(499-501q\). El denominador y \(\alpha\) son positivos. La realización dimensional y el cociente son los de las secciones de acción previamente construidas. La década \(-34\) y la unidad \(\mathcal U_S\) conservan su derivación determinantal anterior.
\end{proof}

Esta fórmula explicita la sección de Planck dentro del contraángulo. En términos de las magnitudes de acción, la relación angular es
\[
 C^*=2\left(\frac{\hbar_{\rm ret}}{10^{-34}\mathcal U_S}+\alpha\right).
\]
La división por la base dimensional extrae el coeficiente adimensional; no suma una magnitud de acción a una constante sin unidades.

Para la orientación angular \(\varepsilon\in\{-1,+1\}\), definimos
\[
 C^*_{\varepsilon}=2\varepsilon(\eta_{\rm ret}+\alpha),\quad
 \xi_\varepsilon=C^*_{\varepsilon}/A,\quad
 q_\varepsilon=\frac{1-\xi_\varepsilon}{1+\xi_\varepsilon}>0.
\]
El recuperador sobre ambas coordenatizaciones es
\begin{equation}
 \eta_{\rm ret}=\alpha\left(500\varepsilon
        \frac{1-q_\varepsilon}{1+q_\varepsilon}-1\right).
 \label{rec:eq:eta-orientada}
\end{equation}

\begin{corollary}[Invariancia del recuperador conjugado]
La transformación \((\varepsilon,q)\mapsto(-\varepsilon,q^{-1})\) conserva \(\eta_{\rm ret}\), \(\hbar_{\rm ret}\) y \(R_{\rm act}\).
\end{corollary}
\begin{proof}
La razón \((1-q)/(1+q)\) cambia de signo al invertir \(q\). El producto por \(\varepsilon\) permanece invariante. Se aplica \eqref{rec:eq:eta-orientada} y después las representaciones de acción.
\end{proof}

Si se descarta la orientación y se mantiene indebidamente \(\varepsilon=+1\), la misma expresión produce \(\eta_{\rm ret}(q^{-1})=-\eta_{\rm ret}(q)-2\alpha\). El registro distingue, por tanto, dos operaciones algebraicamente distintas. La orientación angular \(\varepsilon\) tampoco se identifica sin más con el signo simpléctico \(\sigma\): el intercambio de ejes es antisimpléctico, mientras el giro que los intercambia preserva la forma de área.

\subsection{Entrelazamiento de fase y anisotropía}

Sean
\[
 S_\eta=\operatorname{diag}(e^{\eta/2},e^{-\eta/2}),\quad
 J_0=\begin{pmatrix}0&-1\\1&0\end{pmatrix},\quad
 J_\eta=S_\eta J_0S_\eta^{-1}
       =\begin{pmatrix}0&-e^\eta\\e^{-\eta}&0\end{pmatrix}.
\]
Se tienen \(\det S_\eta=1\), \(J_\eta^2=-I\) y
\(C_\eta(\theta)=e^{\theta J_\eta}=S_\eta e^{\theta J_0}S_\eta^{-1}\).
La ley \(S_{\eta+\delta}=S_\delta S_\eta\) prueba el entrelazamiento exacto
\begin{equation}
 C_{\eta+\delta}(\theta)S_\delta=S_\delta C_\eta(\theta).
 \label{rec:eq:entrelazamiento}
\end{equation}
La transformación preserva \(dQ\wedge dP\) y lleva la evolución de fase a la nueva elipse. Con la normalización de \eqref{eq:ii-elipse-area-fase}, la acción barrida es \(\hbar\theta\), o \(\sigma\hbar\theta\) en la hoja orientada. La fase \(\theta\), la anisotropía \(\eta_{\rm el}\) y el parámetro \(t\) del flujo de \(K\) conservan así funciones diferentes, relacionadas por operadores explícitos.
